\documentclass[12pt]{article}
\usepackage[margin=1in]{geometry}
\usepackage{amsmath}
\usepackage{mathtools}
\usepackage{amssymb}
\usepackage{lmodern}
\usepackage{cancel}
\usepackage{tikz}
\usepackage{graphicx}
\setlength{\parskip}{0.3em}
\setlength{\parindent}{0pt}
% 使用系统字体，需要 XeLaTeX 或 LuaLaTeX 编译
\usepackage{fontspec}
\setmainfont{Times New Roman}
\usepackage{xcolor}
\usepackage{float}
\usepackage{caption}
\usepackage{indentfirst}
\begin{document}
\title{Answers From Chatgpt}
\author{Su Yingze}
\date{\today}
\maketitle
\section{Explanation of the SU(2) Global Anomaly Remark}

We explain the logic of the remark step by step.

\subsection{SU(2) Gauge Theory in Four Dimensions}

The weak interaction is described by an $SU(2)$ gauge theory. Therefore, on our
four-dimensional spacetime $M$ there is an $SU(2)$ principal bundle
\[
P \to M.
\]
On a local patch we may treat spacetime as $\mathbb{R}^4$, and a gauge
transformation is a map
\[
g : \mathbb{R}^4 \to SU(2).
\]

\subsection{Compactification to $S^4$ and Large Gauge Transformations}

Assume the gauge transformation becomes trivial at spatial infinity:
\[
g(x) \to \mathbf{1} \quad \text{as } |x|\to\infty.
\]
This allows us to compactify $\mathbb{R}^4$ by adding a point at infinity:
\[
\mathbb{R}^4 \cup \{\infty\} \cong S^4,
\]
with $g(\infty)=\mathbf{1}$. Hence gauge transformations are classified by maps
\[
g : S^4 \to SU(2) \simeq S^3.
\]
The homotopy classes of such maps are elements of
\[
\pi_4(S^3)\simeq\mathbb{Z}_2=\{0,1\}.
\]
The nontrivial class ($1$) corresponds to a ``large'' gauge transformation.

\subsection{SU(2) Doublet Fermions}

Fermions charged under the weak force transform in the fundamental (doublet)
representation $\mathbf{2}$ of $SU(2)$. A single such fermion field is a section of 
the associated vector bundle
\[
\mathbf{2}\otimes S \longrightarrow M,
\]
where $S$ is the spinor bundle. For our purposes, the only important fact is:

\medskip
\centerline{\textbf{A single SU(2) doublet fermion is one section of a vector bundle.}}
\medskip

\subsection{Witten's SU(2) Global Anomaly}

A key result of Witten (1982) states that under a large gauge transformation
$g:S^4\to S^3$, the partition function of a single $SU(2)$ doublet fermion transforms as
\[
Z_{\text{single doublet}} \;\longrightarrow\; (-1)^s\, Z_{\text{single doublet}},
\]
where $s=[g]\in \pi_4(S^3)\simeq\mathbb{Z}_2$. Thus
\[
s=0:\; Z \text{ invariant},\qquad 
s=1:\; Z\mapsto -Z.
\]
Therefore:

\medskip
\centerline{\textbf{A single SU(2) doublet suffers from a global SU(2) anomaly.}}
\medskip

\subsection{Multiple Doublets and Anomaly Cancellation}

If there are $N$ doublets, the total transformation is
\[
Z \longrightarrow (-1)^{Ns} Z.
\]
The theory is gauge invariant for all gauge transformations if and only if $N$ is even.
In the Standard Model, one generation contains
\[
3 \text{ quark doublets } + 1 \text{ lepton doublet } = 4 \text{ doublets}.
\]
Hence
\[
(-1)^{4s} = +1,
\]
and the anomaly cancels.

This explains, for instance, why one cannot simply change $SU(3)$ color to $SU(4)$:
it would change the number of weak doublets, potentially making $N$ odd and producing
an inconsistency.

\subsection{Summary}

\begin{itemize}
    \item Large $SU(2)$ gauge transformations are classified by $\pi_4(S^3)\simeq\mathbb{Z}_2$.
    \item A single doublet fermion changes sign under the nontrivial transformation.
    \item A consistent theory must have an \emph{even} number of doublets.
    \item The Standard Model satisfies this automatically ($N=4$ per generation).
\end{itemize}

\section{Fiber Choice in the Covering Map $\mathbb{R}\to S^1$}

Consider the covering map
\[
p:\mathbb{R}\longrightarrow S^1,\qquad p(x)=e^{2\pi i x}.
\]
For any point $z=e^{2\pi i\alpha}\in S^1$, the fiber is
\[
p^{-1}(z)=\{x\in\mathbb{R}\mid e^{2\pi i x}=z\}=\alpha+\mathbb{Z},
\]
a translated copy of $\mathbb{Z}$.  
Thus every fiber is homeomorphic to the discrete set $\mathbb{Z}$.

In the long exact sequence of a fibration, one chooses a basepoint $b_0\in S^1$
and defines the fiber as $p^{-1}(b_0)$.  
Although $b_0=1\in S^1$ (the identity element) is the conventional choice,
any $b_0$ works because $S^1$ is path-connected and homotopy groups based
at different points are naturally isomorphic.

\section{Homotopy Groups of the Discrete Space $\mathbb{Z}$}

Although $\mathbb{Z}$ is not a manifold, it is a valid topological space
and its homotopy groups are well-defined.

\subsection{All Continuous Maps $S^n\to\mathbb{Z}$ Are Constant ($n\ge 1$)}

Since $\mathbb{Z}$ is discrete, each singleton $\{k\}\subset\mathbb{Z}$ is open.
A continuous map from a connected space $S^n$ to a discrete space must have
connected image, hence the image consists of only one point.  
Thus any map
\[
f:S^n\to\mathbb{Z},\quad n\ge1
\]
is constant.

\subsection{Computation of $\pi_n(\mathbb{Z})$}

If we fix the basepoint $k_0\in\mathbb{Z}$, the only based map
$S^n\to\mathbb{Z}$ is the constant map $f(x)\equiv k_0$.  
Hence, for $n\ge1$,
\[
\pi_{n}(\mathbb{Z},k_0)=0.
\]

For $n=0$, the path components of $\mathbb{Z}$ are just its points, so
\[
\pi_0(\mathbb{Z})\cong\mathbb{Z}.
\]

\subsection{Summary}

\[
\boxed{
\pi_0(\mathbb{Z})=\mathbb{Z},\qquad
\pi_{n\ge1}(\mathbb{Z})=0.
}
\]

\subsection{Application: Computing $\pi_n(S^1)$ Using the Covering Space}

From the fibration
\[
\mathbb{Z}\longrightarrow \mathbb{R}\longrightarrow S^1,
\]
the long exact sequence gives
\[
\pi_n(S^1)\cong \pi_{n-1}(\mathbb{Z}).
\]
Using our calculation of $\pi_{n-1}(\mathbb{Z})$, we obtain the classical result:
\[
\pi_n(S^1)=
\begin{cases}
\mathbb{Z}, & n=1,\\[4pt]
0, & n\ne 1.
\end{cases}
\]

\section{Understanding why $\partial : \mathbb{Z} \to \mathbb{Z}$ must be multiplication by $2$}

Consider the fibration
\[
SO(2)\;\longrightarrow\; SO(3)\;\longrightarrow\; S^{2}.
\]
The long exact sequence in homotopy (for $n=3$ in the book's notation) contains the segment
\[
\pi_2(SO(2)) \longrightarrow \pi_2(SO(3))
   \longrightarrow \pi_2(S^2) \xrightarrow{\;\partial\;}
   \pi_1(SO(2)) \longrightarrow \pi_1(SO(3)) \longrightarrow \pi_1(S^2).
\]

We substitute the known homotopy groups:
\[
\pi_2(SO(2)) = 0,\quad
\pi_2(S^2) = \mathbb{Z},\quad
\pi_1(SO(2)) = \mathbb{Z},\quad
\pi_1(SO(3)) = \mathbb{Z}_2,\quad
\pi_1(S^2)=0.
\]
Thus the exact sequence becomes
\[
0 \longrightarrow \pi_2(SO(3))
   \longrightarrow \mathbb{Z} \xrightarrow{\;\partial\;}
   \mathbb{Z} \xrightarrow{\;\bmod 2\;}
   \mathbb{Z}_2 \longrightarrow 0.
\]

\subsection{Step 1: Determining the form of $\partial$}

Exactness at $\pi_1(SO(2))=\mathbb{Z}$ means
\[
\mathrm{Im}(\partial)=\ker(\bmod 2).
\]
Since
\[
\ker(\bmod 2)=2\mathbb{Z},
\]
we obtain
\[
\mathrm{Im}(\partial)=2\mathbb{Z}.
\]

Any group homomorphism $\partial : \mathbb{Z}\to \mathbb{Z}$ must have the form
\[
\partial(n)=kn,\qquad k\in\mathbb{Z}.
\]
Its image is $k\mathbb{Z}$.  
To satisfy $\mathrm{Im}(\partial)=2\mathbb{Z}$, the only possibility is
\[
k=\pm2.
\]
Hence
\[
\partial(n)=\pm 2n.
\]
This is the meaning of the sentence in the book:
\begin{quote}
``This forces $\partial : \mathbb{Z}\to\mathbb{Z}$ to be a multiplication by $2$.'' 
\end{quote}

\subsection{Step 2: Computing $\pi_2(SO(3))$}

Exactness at $\pi_2(S^2)=\mathbb{Z}$ implies
\[
\mathrm{Im}\big(\pi_2(SO(3))\longrightarrow\mathbb{Z}\big)
   = \ker(\partial).
\]
Since $\partial(n)=\pm 2n$, its kernel is
\[
\ker(\partial)=\{0\}.
\]
Thus
\[
\mathrm{Im}\big(\pi_2(SO(3))\to\mathbb{Z}\big)=0.
\]
But the map $0\to\pi_2(SO(3))$ at the beginning of the sequence is injective, so the only way
$\pi_2(SO(3))$ can map to $0$ inside $\mathbb{Z}$ is if
\[
\pi_2(SO(3))=0.
\]

So the conclusion is:
\[
\boxed{\pi_2(SO(3))=0.}
\]




\end{document}